% paper/sections/11_full_algorithm.tex
% §11: Navier--Stokes 方程式の計算スキーム全体像

\section{Navier--Stokes 統合更新手順：界面面幾何で閉じる 1 タイムステップ更新}
\label{sec:algorithm}

第~\ref{sec:CCD}章から第~\ref{sec:ppe_solve}章まで，CCD 空間離散化，
CLS 界面追跡，時間積分，圧力連成，および Poisson 方程式の離散解法を個別に定式化した．
本章ではそれらを単なる部品表として並べず，
界面・運動量・圧力反力が同じ面幾何を参照する
1 つの離散 Navier--Stokes 更新として接続する．
読者が最初に押さえるべき点は，
本章の 1 タイムステップを「式を上から順に実行する手続き」としてではなく，
\textbf{界面台帳，運動量台帳，圧力・仕事台帳を同じ面で閉じる操作}
として読むことである．
時刻 $t^n$ から $t^{n+1}$ への式は，この 3 つの台帳が
どの順で所有者を変え，どの量を派生させ，どの面データを次段へ渡すかを
固定するために書く．

\paragraph{本章の位置付け．}
本章は新しい実験結果を追加する章ではなく，
第II部で導入した数値要素を 1 つの Navier--Stokes 計算手順として接続する章である．
個別構成要素の精度は第~\ref{sec:numerical_integrity}章で，
統合後の物理整合性と限界は第~\ref{sec:verification}章で検証する．
したがって本章の役割は，
各演算子の名前を列挙することではなく，
どの未知量が所有者で，どの量がそこから派生した表示量であり，
どの面データが次の演算へ渡されるのかを明確にすることである．
この対応を固定してはじめて，式・演算子・検証章の関係を一続きに読める．

\paragraph{本章の構成．}
まず Navier--Stokes 方程式の各項と離散演算子の対応を示し
（§~\ref{sec:algo_operator_map}），続いて 7 段の更新手順を記述する
（§~\ref{sec:algo_flow}）．初期整合化，ブートストラップ，
タイムステップ制御は §~\ref{sec:bootstrap}，
§~\ref{sec:algo_timestep} にまとめる．
最後に，標準経路の検証済み構成要素を
鋭い界面極限へ集約する発展的構成として，
純 FCCD DNS（§~\ref{sec:pure_fccd_dns}）の位置付けと検証範囲を述べる．
さらに，幾何セル分率を主状態とし，アクティブ集合だけを更新する
アクティブ幾何毛管分解（§~\ref{sec:ao_fast_state_space}）を，
標準経路へ混ぜない発展的毛管閉包として定義する．

\begin{tcolorbox}[defbox, title={11章の理解モデル：3つの台帳を同じ面で閉じる}]
本章で理解すべき核心は，数式そのものの多さではなく，
\textbf{同じ界面面幾何を共有しない更新は二相流の力学として閉じない}
という一点である．
そこで 7 段更新を次の 3 つの台帳として読む．
\begin{enumerate}[noitemsep,leftmargin=2em]
  \item \textbf{界面台帳}（第1--4段）：
        界面を運び，必要なら再初期化し，
        $\psi,\phi,\rho,\mu,\kappa_{lg}$ を同じ界面形状から読める状態へそろえる．
        ここで作った界面面データが，後段の圧力ジャンプと表面張力の基準面になる．
  \item \textbf{運動量台帳}（第5段）：
        圧力で非圧縮性を回収する前に，移流・粘性・重力残差だけで速度を仮に進める．
        この段階の速度 $\bu^*$ は完成品ではなく，圧力に渡す「まだ発散を含む候補」である．
  \item \textbf{圧力・仕事台帳}（第6--7段）：
        圧力増分が，発散の除去，界面圧力ジャンプ，表面張力仕事を同じ面空間で回収する．
        最後に得るべき出力は節点速度だけでなく，次の界面輸送がそのまま使える射影後の面速度である．
\end{enumerate}
したがって各式を読むときの問いは，
「どの物理量を更新しているか」だけでは不十分である．
「その量の所有者は何か」
「それはどの界面面データから派生したか」
「次段へ渡す面速度・圧力ジャンプ・毛管補正量は同じ面を参照しているか」
を追うと，本章全体の構造が見える．
\end{tcolorbox}

\paragraph{本章の式の追い方．}
7 段更新の各式は，値を計算する行であると同時に，
所有者を切り替える行である．
界面段では $\psi$ または $q$ が所有者で，$\phi,\rho,\mu,\kappa_{lg}$ はそこから派生する．
予測段では $\bu^*$ が圧力に渡す候補になり，
投影段では PPE が候補速度を許容面速度へ戻す．
最後に次のステップへ渡すべき主出力は，節点速度だけでなく，
次の界面輸送が直接使う射影後の面速度である．
この所有者の移動を追えば，長いアルゴリズム表は
「どのデータが次段の根拠になるか」を示す流れ図として読める．

% ═════════════════════════════════════════════════════════════════════════════
\subsection{演算子マッピングと全体概観}
\label{sec:algo_operator_map}
% ═════════════════════════════════════════════════════════════════════════════

図~\ref{fig:ns_operator_map} に，NS 方程式の各項と本手法で用いる離散演算子の対応を示す．
標準経路は圧力ジャンプ型の毛管閉包を用いるが，
図では，幾何セル分率から $\Gamma_h$ と派生ゲージ $\phi,\psi$ を作り直し，
その幾何から表面張力項を読むアクティブ幾何経路の採用境界も併記する．
コロケート格子（セル中心）上で圧力・速度・界面変数すべてが同一点に配置され，
圧力・速度のデカップリングは，圧力場を散逸フィルタで平滑化するのではなく，
FCCD/面随伴面フラックスと同じ評価位置の補正段階で抑制する．
本章で重要なのは，個々の演算子の次数だけではない．
界面を横断する圧力勾配，表面張力，HFE 延長，補正段階のジャンプ補正は，
同じ界面面データを共有したときにのみ Balanced--Force の釣合いを保つ．
したがって二相標準経路の圧力閉包は，§~\ref{sec:pressure_accuracy_summary} の
分相 PPE + HFE + DC + 大域ゲージ・再投影分離である．
一括 smoothed-Heaviside PPE は単相・低密度比の一流体対照モデルとして残すが，
二相圧力ジャンプを閉じる並列の標準経路とはしない．

%%% ── 図：NS各項と離散演算子の対応 ─────────────────────────────────────────
\begin{figure}[htbp]
\centering
\begin{tikzpicture}[
  node distance=4mm and 4mm,
  nsterm/.style={draw=navy!70, fill=lightblue, rounded corners=4pt,
                 font=\small\bfseries, inner sep=5pt,
                 minimum width=26mm, minimum height=15mm,
                 text width=23mm, align=center},
  method/.style={draw=#1!70, fill=#1!10, rounded corners=4pt,
                 font=\scriptsize\bfseries, inner sep=5pt,
                 minimum width=26mm, minimum height=15mm,
                 text width=23mm, align=center},
  method/.default=navy,
  gate/.style={draw=darkgray!70, fill=darkgray!7, rounded corners=4pt,
               font=\scriptsize\bfseries, inner sep=5pt,
               minimum width=26mm, minimum height=15mm,
               text width=23mm, align=center},
  subarr/.style={-Stealth, semithick, darkgray!60},
  gatearr/.style={-Stealth, semithick, darkgray!70, dotted},
]

\node[font=\small\bfseries, text=navy] (hd1) at (0,0)
  {NS 方程式の各項と離散演算子の対応};

%% ── NS 5項ボックス ─────────────────────────────────────────────────
\node[nsterm, below=7mm of hd1, xshift=-60mm] (T1) {対流項};
\node[nsterm, right=4mm of T1]  (T2) {圧力項};
\node[nsterm, right=4mm of T2]  (T3) {粘性項};
\node[nsterm, right=4mm of T3]  (T4) {重力項};
\node[nsterm, right=4mm of T4]  (T5) {表面張力項};

%% ── 各離散演算子ボックス ─────────────────────────────────────────────
\node[method=orange2, below=5mm of T1] (S1)
  {移流離散\\[1pt]\footnotesize（UCCD6/FCCD）};
\node[method=midblue, below=5mm of T2] (S2)
  {圧力閉包\\[1pt]\footnotesize（分相 PPE）};
\node[method=teal2, below=5mm of T3] (S3)
  {粘性閉包\\[1pt]\footnotesize（全応力）};
\node[method=darkgray, below=5mm of T4] (S4)
  {体積力\\[1pt]\footnotesize（釣合い浮力）};
\node[method=red2, below=5mm of T5] (S5)
  {標準毛管閉包\\[1pt]\footnotesize（圧力ジャンプ）};
\node[gate, below=4mm of S5] (S5b)
  {発展的毛管閉包\\[1pt]\footnotesize（$q_C\!\to\!\Gamma_h\!\to\!\phi,\psi$）};

\foreach \t/\s in {T1/S1, T2/S2, T3/S3, T4/S4, T5/S5}
  \draw[subarr] (\t) -- (\s);
\draw[gatearr] (S5.south) -- (S5b.north);

\end{tikzpicture}
\caption{NS 方程式の各項と対応する離散演算子の概要}
\label{fig:ns_operator_map}
\PaperCaptionNote{速度主変数経路の対流項は IMEX-BDF2 の EXT2 外挿で評価し，領域内部は UCCD6（BF 面整合を要する一流体対照では FCCD 面構成）で扱う．保存形共通フラックス経路では，CLS $\psi$ の FCCD 面フラックス台帳から密度と運動量も同時に更新する．粘性項は陰的 BDF2 の全応力 Helmholtz 系を欠陥補正（DC）で閉じる．二相標準圧力項は分相 PPE + HFE + DC + 勾配演算子 $\mathcal{G}_\Gamma$ による投影法~\cite{Chorin1968,vanKan1986}で扱い，DC は作用中の PPE 残差を下げる補正だけを受理し，圧力場には DCCD 後置フィルタを適用しない．一括 smoothed-Heaviside PPE と CSF 体積力は一流体対照経路であり，圧力ジャンプ型二相閉包の代替ではない．重力は釣合い浮力分解の面残差として予測段階に入れる．分相圧力ジャンプ IPC では，履歴へ保存するのは滑らかな圧力座標であり，界面ジャンプと幾何毛管反力は現在の切断面法則で再構成する．アクティブ幾何経路は，幾何セル分率 $q_C$ から互換界面と派生ゲージ $\phi,\psi$ を再構成し，毛管共変量 $r_\sigma$，圧力反力部分空間 $R_p$，移動格子の投影済み面共鎖移送を同時に満たすとき標準毛管閉包として用いる．壁付き計算では $F_w=\ker C_w$ を面速度状態空間とし，圧力補正後の壁修復は標準経路にしない．射影後の許容面速度を CLS フラックスへ直接渡す．}
\end{figure}

本章の標準経路では，DCCD を章全体の安定化調整パラメータとして扱わない．
DCCD は §~\ref{sec:dccd_derivation} で定義した節点中心 CCD の外殻後置フィルタであり，
CLS 移流は §~\ref{sec:advection_motivation} の FCCD 面フラックス，
運動量の領域内部は §~\ref{sec:uccd6_def} の UCCD6，
圧力・毛管ジャンプを含む面上補正は §~\ref{sec:balanced_force}--§~\ref{sec:ppe_solve} の
同一評価位置条件に従う．
高密度比の移動界面では，これに加えて保存形共通フラックス分岐を用いる．
この分岐では，CLS の段階面フラックス台帳から
$q,m,\bm{p}_m$ を同時に更新し，運動量予測子では速度主変数対流を再評価しない．
したがって標準 UCCD6 対流は滑らかな速度主変数経路の部品であり，
保存形共通フラックス経路の対流を二重に置き換える部品ではない．

\paragraph{アクティブ幾何毛管分解の位置付け．}
標準圧力ジャンプ経路とは別に，表面エネルギー仮想仕事を
アクティブ界面幾何から直接構成する発展的経路を
本稿ではアクティブ幾何毛管分解と呼ぶ．
ただし入口は毛管力の後処理ではなく，輸送後の $q_C$ と
ゲージ予測子から互換界面を再構成し，そこから $\phi,\psi$ を
派生させる界面状態再構成である．
この経路は，幾何セル分率 $q_C$，互換射影，圧力反力部分空間 $R_p$，
近似誤差境界，移動格子での投影済み面共鎖移送，
および滑らかな圧力履歴座標が同時に成立する場合だけ
標準経路の候補になる．
いずれかの条件が不成立なら標準圧力ジャンプ経路へ自動で置き換えず，
未受理の発展的閉包として停止する．
したがって，第~\ref{sec:numerical_integrity}章の U12 と
第~\ref{sec:verification}章の V11 は，この経路が物理ベンチマークへ
混入しないことを確認する事前ゲートであり，
受理境界の詳細は第~\ref{sec:ao_fast_state_space}節でまとめる．

表~\ref{tab:algo_overview} に 7 段の読み取り方を示す．
この表で見るべきものは手法名そのものではなく，
各段がどの台帳を閉じ，どの出力を次段へ渡すかである．
順序の要点は，界面輸送と幾何更新を先に閉じ，
その後の予測段階，PPE，補正段階が同一の
$\psi,\phi,\rho,\mu,\kappa_{lg}$ と界面面データを参照することである．

\begin{table}[htbp]
\centering
\caption{1 タイムステップの更新手順と読み取り方}
\label{tab:algo_overview}
\PaperCaptionNote{$t^n \to t^{n+1}$ の更新手順を，台帳の役割と次段へ渡す出力に分けて示す．}
\footnotesize
\begin{tabular}{clp{0.30\linewidth}p{0.31\linewidth}l}
\hline
\textbf{段} & 役割 & 手法 & 読者が追う出力 & 参照 \\
\hline
第1段 & CLS $\psi$ 移流
  & TVD-RK3 + FCCD 面フラックス発散（保存形共通フラックスでは同じ台帳で $q,m,\bm{p}_m$ を更新）
  & どの射影後面速度で界面を運んだか．共通フラックス分岐では界面・密度・運動量が同じ面フラックス台帳で進む．
  & §~\ref{sec:cls_advection} \\
第2段 & 再初期化
  & Ridge--Eikonal 補助距離再構成 + 鋭い界面体積シフト + 拡散質量プロファイル閉鎖
  & 零点集合，鋭界面体積，拡散 CLS 質量を混同せず閉じたか．共通フラックスで保存的再写像がない場合は受理しない．
  & §~\ref{sec:cls_compression} \\
第3段 & 界面幾何・物性更新
  & $\psi\!\to\!\phi$ ロジット変換 + 線形補間
  & 後段が参照する $\psi,\phi,\rho,\mu$ が同じ界面状態から派生したか．
  & §~\ref{sec:levelset_mapping} \\
第4段 & 曲率計算
  & $\psi$ 直接曲率 + 界面帯制限フィルタ
  & 圧力ジャンプと毛管仕事が読む $\kappa_{lg}$ が，第3段の界面帯と同じ面幾何に属しているか．
  & §~\ref{sec:curvature_direct_psi_cls} \\
第5段 & 運動量予測
  & 速度主変数経路は UCCD6 + EXT2 + 粘性 Helmholtz；保存形共通フラックス経路では第1段で運動量移流済み
  & 圧力補正前の候補速度 $\bu^*$ と，圧力空間に残すべき圧力ジャンプ・静水圧成分を分けたか．
  & §~\ref{sec:time_level1} \\
第6段 & PPE 評価
  & 分相 PPE + HFE + 残差低下 DC + IPC 増分
  & 発散除去，圧力ジャンプ，表面エネルギー仕事を同じ圧力仕事計量と面空間で回収したか．
  & §~\ref{sec:ppe_solve} \\
第7段 & 補正段階
  & IPC 増分更新（同一 $\gamma\Delta t$ 投影；成分補正済み毛管面加速度を保存）
  & 節点速度だけでなく，次の第1段が直接使う射影後面速度を保存したか．
  & §~\ref{sec:ipc_derivation} \\
\hline
\end{tabular}
\end{table}

$\psi$ は保存形レベルセットの主変数であり，
$\phi$ は法線計算，HFE，界面厚補正などで必要となる幾何補助量である．
したがって，単なる $\psi\to\phi\to\psi$ は恒等変換であり，界面プロファイルの
修復には零点集合からの Ridge--Eikonal 距離再構成と
鋭界面相体積を閉じる SDF 一様シフト，
および拡散 CLS 質量を閉じるプロファイル幅補正を挟む．

\paragraph{本章で用いる離散演算子の記法}
\label{box:operator_notation}
\begin{itemize}[noitemsep]
  \item $D_x^{(1)}, D_y^{(1)}$：CCD 法による1次微分 $(\partial/\partial x,\, \partial/\partial y)$
  \item $D_x^{(2)}, D_y^{(2)}$：CCD 法による2次微分 $(\partial^2/\partial x^2,\, \partial^2/\partial y^2)$
  \item $D_{xy}^{(11)}$：CCD 法による混合微分 $\partial^2/\partial x\partial y$—
        $(D_{xy}^{(11)} f)_{i,j} = (D_y^{(1)}[D_x^{(1)} f])_{i,j}$（逐次2段スウィープ，$\Ord{h^6}$）
  \item $\mathcal{A}_{\bu}(\bu)$：運動量移流オペレータ．
          標準の領域内部は UCCD6（§~\ref{sec:uccd6_def}）で評価し，
          BF 面整合を検証する一流体対照系では FCCD 面構成（§~\ref{sec:fccd_advection}）を用いる．
          CCD 中心差分の非保存形
          $u\,D_x^{(1)}q+v\,D_y^{(1)}q$ は滑らかな単相製造解・比較記法に限る．
          CLS $\psi$ 移流には FCCD 面フラックス発散 $\mathrm{Div}_h^{\psi}$ を使用する．
  \item $\mathcal{D}_{\mathrm{d},x}$，$\mathcal{D}_{\mathrm{d},y}$：各速度成分自身に作用する対角粘性項
        （本稿では陰的 BDF2 の Helmholtz 型 DC で陰的に評価；§~\ref{sec:time_viscous}）
  \item $\mathcal{D}_{\mathrm{c},x}$，$\mathcal{D}_{\mathrm{c},y}$：他成分を含む交差粘性項
        （陰的 BDF2 + DC では高次残差内で同時刻評価；§~\ref{sec:time_viscous}）
  \item $\mathrm{Div}_h^{\psi}((P_f\psi)\bu_f)$：FCCD 保存形面フラックス発散（CLS $\psi$ 移流専用．
        射影後の面速度を保存する経路では式~\eqref{eq:projection_native_psi_transport} の面速度 $\bu_f$ を使う；§~\ref{sec:advection_motivation}）
  \item $\mathcal{G}^\text{adj}$：面平均圧力勾配（非一様格子 + 壁面境界条件専用，$\Ord{h^2}$）
          \[
            (\mathcal{G}^\text{adj}_x\, p)_{i,j}
            = \tfrac{1}{2}\!\left[
                \tfrac{p_{i+1,j}-p_{i,j}}{d_f^{x,(i)}}
              + \tfrac{p_{i,j}-p_{i-1,j}}{d_f^{x,(i-1)}}
              \right]
          \]
          $J_f = 1/d_f^{(i)}$ を使い $\mathcal{D}_\text{FV}(\mathcal{G}^\text{adj}) = \mathcal{L}_\text{FV}$ を保証する
          （§~\ref{sec:fvm_ccd_corrector}；均一格子では $\mathcal{G}^\text{adj} = D^{(1)}_\text{CCD}$）．
          圧力ジャンプ使用時は同じ面・切断面位置で
          $\mathcal{G}^{\text{adj}}_\Gamma=\mathcal{G}^{\text{adj}}-B^{\text{adj}}$ とし，
          $B^{\text{adj}}$ は式~\eqref{eq:cut_face_young_laplace_jump} の
          切断面 Young--Laplace ジャンプと局所物理面距離 $H_f$ で構成する．
  \item $\mathcal{L}_{\mathrm{CCD}}^{\rho}$：CCD 積の微分則型変密度ラプラシアン
          $= (1/\rho)(D_x^{(2)}p + D_y^{(2)}p)
            - (D_x^{(1)}\rho/\rho^2)D_x^{(1)}p
            - (D_y^{(1)}\rho/\rho^2)D_y^{(1)}p$
          （演算子単体 $\Ord{h^6}$．
          分相 PPE 適用時は，DC 残差受理により液相内部 Dirichlet 評価で
          $\Ord{h^7}$ が回復することが §~\ref{sec:verify_split_ppe_dc} で実証されているが，
          Neumann 境界条件では境界スキームにより $\Ord{h^5}$ が上界となる．
          一括 smoothed-Heaviside 密度場では界面近傍で高次精度が失われ，
          高解像度側で格子収束がプラトー化する；§~\ref{sec:interface_crossing}）
\end{itemize}

% ═════════════════════════════════════════════════════════════════════════════
